% Осень26!B35; нумерация внутри столбца без пустых строк.
% Источник: лекция 7, §6 (строка 171)
\definitionticket{32}{Сходимость почти везде}
\label{definition:32}

\subsection*{Свойства, выполняющиеся почти везде; сходимость почти везде}

\textbf{Мотивировка лектора.} Это «пропущенный параграф» теории --- его место было \textbf{перед} определением интеграла: чтобы говорить «совпадают на множестве меры ноль / отличаются на множестве меры ноль», никакие интегралы не нужны. «Заодно, чтобы два раза не вставать», сюда же добавлена сходимость по мере (§7).

\subsubsection*{Определения}

\textbf{Терминология.} «Почти везде» = «почти всюду» = «при почти всех \(x\)» --- вариации одного термина; общепринятое сокращение: \textbf{п.в.}

\begin{definition}

\emph{свойство, выполняющееся почти везде.} Пусть \((X,\mathcal A,\mu)\) --- пространство с мерой, \(E\in\mathcal A\), и \(W(x)\) --- высказывание о точке \(x\). Говорят, что \(W\) \textbf{верно при почти всех \(x\in E\)}, если существует множество \(E'\in\mathcal A\), \(E'\subset E\), \(\mu(E')=0\), такое что \(W(x)\) верно для всех \(x\in E\setminus E'\).

\end{definition}

\textbf{Главный пример высказывания.} Для последовательности функций \(f_1,f_2,\dots\colon X\to\overline{\mathbb R}\) и функции \(f\colon X\to\overline{\mathbb R}\): \[W(x):\quad \exists\ \lim_{n\to\infty}f_n(x)=f(x)\] (бесконечные пределы не смущают: \(f\) может принимать бесконечные значения). Если \(W(x)\) верно п.в., говорят, что \textbf{\(f_n\) сходится к \(f\) почти везде} (п.в.): \(f_n\to f\) п.в. Всюду-поточечная сходимость --- частный случай, где «нарушение происходит на пустом множестве».

\subsubsection*{Измеримость п.в.-предела и полнота меры}

\textbf{Пункт 1 (измеримость предела, «подломленная теорема»).} Пусть \(\mu\) \textbf{полная}, \(f_n\) измеримы и \(f_n\to f\) п.в. на \(X\). Тогда \(f\) измерима.

\begin{proof}

Пусть \(E\) --- множество нарушения, \(\mu(E)=0\); положим \(\widetilde X:=X\setminus E\) (измеримо). На \(\widetilde X\) предел существует всюду, и по теореме об измеримости пределов \(f\) измерима на \(\widetilde X\). Проверим измеримость на всём \(X\) через множества Лебега: для \(a\in\mathbb R\) \[X(f<a)\ =\ \bigl(X(f<a)\cap\widetilde X\bigr)\ \sqcup\ \bigl(X(f<a)\cap E\bigr).\] Первое слагаемое измеримо (измеримость \(f\) на \(\widetilde X\)); второе --- \textbf{подмножество множества \(E\) меры ноль}, и оно измеримо \textbf{по полноте} меры.

\end{proof}

\textbf{«Ну и обманул, как и обещал» (лектор).} Без полноты второй кусок может быть \textbf{неизмеримым}: \(\sigma\)-алгебра не обязана следовать нашей логике --- «сигма-алгебра может быть такой, что множество имеет меру ноль, а какой-то кусок у него неизмеримый».

\begin{definition}

\emph{полная мера.} Мера \(\mu\) на \(\sigma\)-алгебре \(\mathcal A\) \textbf{полная}, если всякое подмножество множества меры ноль измеримо: \[A\in\mathcal A,\ \mu(A)=0,\ B\subset A\quad\Longrightarrow\quad B\in\mathcal A\ (\text{и тогда }\mu(B)=0).\]

\end{definition}

\textbf{Пункт 2.} Пусть \(\mu\) полная, \(f\) измерима, \(g=f\) п.в. Тогда \(g\) измерима. (Доказательство «абсолютно аналогично»: на части, где равны, выручает измеримость \(f\); на части внутри множества меры ноль --- полнота. \textbf{Упражнение лектора: проговорите сами.})

\textbf{Если \(\mu\) неполная --- «починка».} Пусть \(f_n\) измеримы, \(f_n\to f\) п.в., но \(\mu\) неполная: тогда \(f\) \textbf{не обязана} быть измеримой, но её можно \textbf{изменить на множестве меры ноль так, что она станет измеримой}; иными словами, \(f\) совпадает вне множества меры ноль с некоторой измеримой функцией. \emph{Пояснение лектора:} на множестве нарушения \(E\) зададим \(f\) «предсказуемо» --- например, \(f|_E:=2026\) (или «нулю какого-нибудь естественного значения»). Тогда для любого \(a\) пересечение \(\{f<a\}\cap E\) --- это либо всё \(E\), либо пустое множество (смотря как \(a\) стоит относительно 2026), оба измеримы; вопрос «поведение \(f\) на \(E\)» исчезает.

\subsubsection*{Эквивалентность функций}

\begin{definition}

Измеримые \(f,g\colon X\to\overline{\mathbb R}\) \textbf{эквивалентны} (в смысле теории меры), \(f\sim g\), если \(f=g\) п.в.

\end{definition}

\emph{Комментарий лектора:} отношение эквивалентности --- общематематическая конструкция, «в каждом предмете свои эквивалентности», обычно со «скромным» значком; аналитики «поступили по-наглому и забрали себе обычный значок \(\sim\)». Для измеримых функций определение корректно: измеримые функции образуют алгебру (суммы, разности измеримых измеримы), множества \(\{f-g>0\}\), \(\{f-g<0\}\) измеримы.

\subsubsection*{Счётное семейство п.в.-свойств}

\textbf{Фундаментальное свойство слов «почти всюду».} Пусть \(W_n(x)\), \(n\in\mathbb N\), --- высказывания о точке, и каждое \(W_n\) верно п.в. на \(X\). Тогда высказывание «\(\forall n\ W_n(x)\)» верно п.в. на \(X\).

\begin{proof}

Пусть \(E_n:=\{x:\ W_n(x)\ \text{неверно}\}\), \(\mu(E_n)=0\). Нарушение «\(\forall n\,W_n\)» возможно лишь на \(\bigcup_{n=1}^{\infty}E_n\), а счётное объединение множеств меры ноль имеет меру ноль (счётная полуаддитивность меры).

\end{proof}

\emph{Пояснение:} при каждом \(n\) высказывание «кое-где нарушается», но все нарушения сразу «садятся» на одно общее множество меры ноль. (Обратная импликация тривиальна: где верны все --- верно и каждое.) «Можно так здорово манипулировать этими словами „почти всюду``» --- но только со \textbf{счётными} семействами.
